\documentclass[11pt]{article}

\usepackage[T1]{fontenc}
\usepackage[utf8]{inputenc}
\usepackage[paper=a4paper,left=25mm,right=25mm,top=25mm,bottom=25mm]{geometry}
\usepackage[english, ngerman]{babel}
\usepackage[autostyle=true,german=quotes]{csquotes}
\usepackage{setspace}
\usepackage[expansion=false]{microtype}
\usepackage{booktabs}
\usepackage{tabularx}
\usepackage{enumitem}
\usepackage{xcolor}
\usepackage{hyperref}

\hypersetup{
    colorlinks,
    citecolor   = black,
    filecolor   = black,
    linkcolor   = black,
    urlcolor    = black,
    pdftitle    = {Seminararbeit SYCL: Zwischenstand und Erkenntnisse},
    pdfauthor   = {Erlind Sejdiu},
    pdfcreator  = {pdflatex},
    pdfproducer = {LaTeX with hyperref}
}
\setstretch{1.25}

% Markierung für Stellen, die noch ergänzt/geprüft werden müssen
\newcommand{\offen}[1]{\textcolor{red}{[#1]}}
% Kasten-artige Hervorhebung einer Kernaussage
\newcommand{\kernaussage}[1]{\par\smallskip\noindent\textbf{Kernaussage:} \emph{#1}\par\smallskip}

\begin{document}

\begin{center}
    {\LARGE\bfseries Seminararbeit SYCL\\[0.3em]
    Zwischenstand und bisherige Erkenntnisse\par}
    \vspace{0.8em}
    Erlind Sejdiu (106357) \quad·\quad Stand: 22.\,10.\,2026
\end{center}

\tableofcontents

\section{Worum geht es?}

Moderne Rechner enthalten neben der CPU oft weitere Recheneinheiten,
vor allem Grafikkarten (GPUs). Eine GPU kann sehr viele gleichartige
Rechnungen gleichzeitig ausführen und ist dadurch bei passenden Aufgaben um
Größenordnungen schneller als eine CPU.

Das Problem: Jeder Hersteller hat seine eigene Programmierschnittstelle.
Code für NVIDIA-GPUs wird typischerweise in CUDA geschrieben und läuft nur
auf NVIDIA-Hardware. Für AMD- oder Intel-GPUs oder für die CPU müsste man
das Programm neu schreiben.

\textbf{SYCL} (gesprochen \enquote{sickle}) ist ein offener Standard der
Khronos Group, die auch OpenGL und Vulkan verantwortet. Die Idee: Man
schreibt \emph{einmal} normales C++, und derselbe Quellcode läuft auf
CPUs und auf GPUs verschiedener Hersteller. Es gibt mehrere Compiler, die
SYCL umsetzen, u.\,a.\ \emph{DPC++} von Intel und \emph{AdaptiveCpp}
(Universität Heidelberg).

Die Leitfrage meiner Arbeit ist:
\begin{quote}
\emph{Wie aufwendig ist es, ein bestehendes C++-Programm auf SYCL
umzustellen, was bringt es an Leistung, und hält das Versprechen
\enquote{einmal schreiben, überall ausführen}?}
\end{quote}

\section{Die wichtigsten Begriffe in einfachen Worten}

\begin{description}[style=nextline, leftmargin=0.5cm]
    \item[Host und Gerät (Device)]
    Der \emph{Host} ist die CPU, auf der das Hauptprogramm läuft. Ein
    \emph{Gerät} ist die Hardware, auf der die eigentliche Rechnung läuft,
    z.\,B.\ eine GPU. Auch die CPU selbst kann als Gerät dienen.
    \item[Queue (Warteschlange)]
    Über eine Queue schickt der Host Arbeitsaufträge an genau ein Gerät.
    Welches Gerät das ist, kann man beim Programmstart über eine
    Umgebungsvariable festlegen, ohne den Code zu ändern.
    \item[Kernel]
    Das Stück Code, das auf dem Gerät ausgeführt wird. In SYCL ist das eine
    gewöhnliche C++-Funktion (Lambda).
    \item[Work-Item]
    Eine einzelne Ausführung des Kernels, vergleichbar mit einem sehr
    leichtgewichtigen Thread. Bei der Matrixmultiplikation berechnet ein
    Work-Item genau einen Eintrag der Ergebnismatrix. Bei einer
    $2048 \times 2048$-Matrix sind das rund 4 Millionen Work-Items.
    \item[Work-Group]
    Eine Gruppe von Work-Items, die zusammenarbeiten, sich einen kleinen,
    schnellen Zwischenspeicher teilen (\emph{Local Memory}) und aufeinander
    warten können (\emph{Barriere}).
    \item[Unified Shared Memory (USM)]
    Die Art, wie Speicher zwischen Host und Gerät geteilt wird. Entweder
    \emph{shared} (SYCL verschiebt die Daten automatisch) oder
    \emph{device} (der Programmierer kopiert die Daten selbst).
    \item[Backend]
    Die herstellerspezifische Schicht, auf die SYCL intern zurückgreift,
    z.\,B.\ CUDA für NVIDIA oder OpenCL für die CPU. Der Programmierer sieht
    sie nicht.
    \item[GFLOPS]
    Milliarden Gleitkommarechnungen pro Sekunde, das übliche Maß für
    Rechenleistung. Höher ist besser.
\end{description}

\section{Vorgehen}

Als Beispiel dient die Multiplikation zweier quadratischer Matrizen
($C = A \cdot B$). Sie ist leicht zu verstehen, braucht viel Rechenzeit
($2n^3$ Operationen) und ist ein Standardbeispiel für GPU-Programmierung.

Ein einfaches C++-Programm wird in Stufen auf SYCL umgebaut. Jede Stufe
führt genau ein neues Konzept ein. Alle Stufen messen auf die gleiche Weise
und prüfen ihr Ergebnis gegen eine Referenzrechnung.

\begin{table}[h]
\centering
\begin{tabularx}{\textwidth}{@{}c>{\raggedright\arraybackslash}X>{\raggedright\arraybackslash}X@{}}
\toprule
\textbf{Stufe} & \textbf{Was passiert} & \textbf{Neues Konzept} \\
\midrule
0 & Ausgangsprogramm: drei verschachtelte Schleifen auf einem CPU-Kern & -- \\
1 & Die äußeren Schleifen werden durch einen parallelen Kernel ersetzt & Queue, Kernel, USM \emph{shared} \\
2 & Daten werden explizit auf das Gerät und zurück kopiert & USM \emph{device}, Kopieren \\
3 & Work-Groups laden Teilblöcke (\enquote{Kacheln}) in den schnellen Local Memory & Work-Group, Local Memory, Barriere \\
4 & Kein neuer Code: Stufe 1--3 auf CPU und GPU ausführen & Geräteauswahl \\
\bottomrule
\end{tabularx}
\caption{Die Stufen des Beispielprojekts}
\end{table}

\paragraph{Messumgebung.}
AMD Ryzen 5 5600X (6 Kerne) und NVIDIA GeForce RTX 5060 Ti. Compiler:
DPC++ 7.2 (Intel LLVM). Die CPU wird über das OpenCL-Backend angesprochen,
die GPU über das CUDA-Backend. Gemessen wird jeweils der Median aus fünf
Wiederholungen, der allererste Lauf wird getrennt erfasst.

\section{Erkenntnisse}

\subsection{Derselbe Code läuft auf CPU und GPU}

Alle Stufen liefern auf beiden Geräten korrekte Ergebnisse. Um zwischen
CPU und GPU zu wechseln, wird keine einzige Codezeile geändert, sondern
nur eine Umgebungsvariable beim Start gesetzt.

Eine Einschränkung gibt es: Das Programm muss beim Kompilieren für alle
gewünschten Geräte übersetzt werden. Ein nur für NVIDIA übersetztes
Programm bricht auf der CPU mit der Meldung ab, es gebe \enquote{kein
Abbild für das gewählte Gerät}. Das ist ein Compiler-Schalter, keine
Codeänderung, aber man muss es wissen.

\kernaussage{Das Versprechen \enquote{einmal schreiben, überall ausführen} hält
im getesteten Rahmen.}

\subsection{Der Einstieg ist klein, der Effekt groß}

Der Schritt von Stufe 0 zu Stufe 1 betrifft im Wesentlichen die zwei
äußeren Schleifen, die durch einen \texttt{parallel\_for}-Aufruf ersetzt
werden. Die eigentliche Rechnung bleibt gleich.

\begin{table}[h]
\centering
\begin{tabular}{@{}lrr@{}}
\toprule
\textbf{Variante ($n = 1024$)} & \textbf{Zeit} & \textbf{GFLOPS} \\
\midrule
Stufe 0, ein CPU-Kern & 4484\,ms & 0{,}5 \\
Stufe 1, CPU (alle Kerne, SYCL) & 87\,ms & 25 \\
Stufe 1, GPU & 1{,}9\,ms & 1105 \\
Stufe 3, GPU & 1{,}5\,ms & 1395 \\
\bottomrule
\end{tabular}
\caption{Rechenzeit ohne Datentransfer}
\end{table}

Die GPU ist hier über 2000-mal schneller als das Ausgangsprogramm. Diese
Zahl ist aber mit Vorsicht zu lesen: Stufe 0 ist absichtlich eine
einfache Implementierung auf nur einem Kern. Ein Vergleich mit einer
optimierten CPU-Bibliothek würde deutlich kleiner ausfallen.

\kernaussage{Mit wenig Umbauaufwand ist bereits ein Großteil der möglichen
GPU-Leistung erreichbar.}

\subsection{Der erste Lauf ist deutlich langsamer}

Beim ersten Aufruf übersetzt die Laufzeitumgebung den Kernel noch
passend für das konkrete Gerät (\emph{Just-in-Time-Kompilierung}) und
initialisiert das Gerät. Bei $n = 256$ dauert Stufe 1 auf der GPU im
ersten Lauf 0{,}75\,ms statt 0{,}04\,ms (etwa 20-mal länger), auf der CPU
118\,ms statt 0{,}4\,ms.

\kernaussage{Für kurze Programme, die nur einmal rechnen, kann diese Anlaufzeit
den Vorteil der GPU aufheben. Bei Messungen muss der erste Lauf getrennt
betrachtet werden.}

\subsection{Bei kleinen Aufgaben dominiert der Datentransport}

Die GPU hat eigenen Speicher. Die Matrizen müssen erst dorthin und das
Ergebnis zurück kopiert werden (Stufe 2, GPU):

\begin{table}[h]
\centering
\begin{tabular}{@{}rrrr@{}}
\toprule
$n$ & \textbf{Rechnung} & \textbf{Kopieren} & \textbf{Anteil Kopieren} \\
\midrule
256 & 0{,}04\,ms & 0{,}17\,ms & 82\,\% \\
1024 & 1{,}97\,ms & 1{,}34\,ms & 40\,\% \\
2048 & 15{,}87\,ms & 3{,}93\,ms & 20\,\% \\
\bottomrule
\end{tabular}
\caption{Rechnung und Datentransfer auf der GPU}
\end{table}

Der Rechenaufwand wächst mit $n^3$, die Datenmenge nur mit $n^2$. Je
größer die Aufgabe, desto mehr lohnt sich die GPU.

\kernaussage{Ob sich eine GPU lohnt, hängt von der Problemgröße ab. Bei kleinen
Matrizen kostet das Kopieren mehr als die Rechnung.}

\subsection{Automatische und manuelle Speicherverwaltung sind gleich schnell}

Stufe 1 (SYCL verschiebt die Daten selbst) und Stufe 2 (Daten werden
manuell kopiert) sind in der reinen Rechenzeit praktisch gleich schnell,
z.\,B.\ 1{,}94\,ms gegenüber 1{,}97\,ms auf der GPU bei $n = 1024$. Der
Unterschied liegt darin, \emph{wann} die Kopierkosten anfallen: Bei Stufe
1 werden sie beim ersten Zugriff versteckt mitbezahlt, bei Stufe 2 sind
sie sichtbar und planbar.

\kernaussage{Die automatische Variante ist für den Einstieg ausreichend. Die
manuelle Variante macht die Kosten sichtbar und kontrollierbar.}

\subsection{Hardwarenahe Optimierung wirkt je nach Gerät sehr unterschiedlich}

In Stufe 3 arbeiten jeweils $16 \times 16$ Work-Items als Gruppe
zusammen. Sie laden gemeinsam einen Teilblock der Matrizen in einen
kleinen, schnellen Speicher und rechnen von dort weiter. Dadurch muss
jeder Wert seltener aus dem langsamen Hauptspeicher gelesen werden.

\begin{table}[h]
\centering
\begin{tabular}{@{}lrrr@{}}
\toprule
\textbf{Gerät} & \textbf{Stufe 2} & \textbf{Stufe 3} & \textbf{Gewinn} \\
\midrule
GPU, $n = 1024$ & 1{,}97\,ms & 1{,}54\,ms & $+28\,\%$ \\
GPU, $n = 2048$ & 15{,}87\,ms & 12{,}12\,ms & $+31\,\%$ \\
CPU, $n = 1024$ & 83{,}1\,ms & 18{,}4\,ms & $\times 4{,}5$ \\
CPU, $n = 2048$ & 450{,}9\,ms & 129{,}2\,ms & $\times 3{,}5$ \\
\bottomrule
\end{tabular}
\caption{Wirkung der Kachel-Optimierung (Stufe 2 $\rightarrow$ Stufe 3)}
\end{table}

Vermutete Erklärung (nicht gesondert gemessen): Die GPU fängt schon in
Stufe 2 durch eigene Caches und gebündelte Speicherzugriffe viel ab. Die
CPU profitiert dagegen stark davon, dass die Teilblöcke in ihren Cache
passen. Bei sehr kleinen Matrizen ($n = 256$) bringt die Optimierung auf
der GPU nichts, dort ist Stufe 3 sogar etwas langsamer.

\kernaussage{Dieselbe Optimierung bringt auf einem Gerät 30\,\%, auf dem anderen
ein Vielfaches.}

\subsection{Portabilität ist nicht gleich Performance-Portabilität}

Hier liegt die bisher wichtigste Erkenntnis. Man muss zwei Begriffe
unterscheiden:

\begin{itemize}
    \item \textbf{Portabilität:} Der Code läuft auf jedem Gerät und
    rechnet richtig.
    \item \textbf{Performance-Portabilität:} Der Code nutzt auf jedem Gerät
    auch einen ähnlich großen Teil der möglichen Leistung, ohne
    gerätespezifische Anpassung.
\end{itemize}

Um das zu prüfen, wurde Stufe 3 mit verschiedenen Kachelgrößen $T$ gebaut
(eine Work-Group umfasst $T \times T$ Work-Items) und auf beiden Geräten
gemessen. Die Prozentangabe zeigt, wie nah die Variante an der besten
Kachelgröße für \emph{dieses} Gerät liegt:

\begin{table}[h]
\centering
\begin{tabular}{@{}rrrr@{}}
\toprule
$T$ & \textbf{GPU} & \textbf{CPU} & \textbf{PP} \\
\midrule
8  & 1135\,GFLOPS (83\,\%) & 71\,GFLOPS (45\,\%) & 0{,}59 \\
16 & \textbf{1368}\,GFLOPS (100\,\%) & 137\,GFLOPS (88\,\%) & 0{,}93 \\
32 & 1354\,GFLOPS (99\,\%) & \textbf{157}\,GFLOPS (100\,\%) & 0{,}99 \\
64 & \emph{startet nicht} & 142\,GFLOPS (91\,\%) & 0 \\
\bottomrule
\end{tabular}
\caption{Kachelgröße und Leistung, $n = 2048$}
\end{table}

Drei Beobachtungen:
\begin{enumerate}
    \item \textbf{Die beste Einstellung unterscheidet sich je nach
    Gerät:} Auf der GPU ist $T = 16$ am schnellsten, auf der CPU $T = 32$.
    Bei $n = 1024$ erreicht die CPU mit $T = 16$ nur 74\,\% ihres
    besten Werts.
    \item \textbf{Eine Einstellung kann ein Gerät ganz ausschließen:} Bei
    $T = 64$ umfasst eine Work-Group 4096 Work-Items. NVIDIA-GPUs erlauben
    höchstens 1024, der Kernel startet dort also gar nicht. Auf der CPU
    läuft dieselbe Variante dagegen gut.
    \item \textbf{Eine schlechte Einstellung trifft die Geräte
    unterschiedlich stark:} $T = 8$ kostet die GPU 17\,\%, die CPU 55\,\%.
\end{enumerate}

Die Spalte \textbf{PP} ist eine in der Forschung gängige Kennzahl für
Performance-Portabilität (Pennycook, Sewall und Lee, \emph{A Metric for
Performance Portability}, 2016). Sie fasst die Prozentwerte aller Geräte
zu einem Wert zwischen 0 und 1 zusammen (harmonischer Mittelwert). Ein
einziges schwaches Gerät zieht den Wert stark nach unten, und läuft der
Code auf einem Gerät gar nicht, ist er 0.

Einschränkung: Die Prozentwerte beziehen sich auf die beste gemessene
Kachelgröße, nicht auf die theoretische Höchstleistung der Hardware.

\kernaussage{SYCL sorgt dafür, dass der Code überall läuft. Damit er überall
schnell läuft, müssen Parameter wie die Kachelgröße weiterhin pro Gerät
abgestimmt werden.}

\subsection{Praktische Hürden}

\begin{itemize}
    \item \textbf{Einrichtung der Werkzeuge:} Die Standardinstallation von
    Intel oneAPI sah auf dem Testrechner nur die CPU. Für die NVIDIA-GPU
    war eine eigene DPC++-Version mit CUDA-Unterstützung nötig. Der Aufwand
    steckt also eher in der Einrichtung als im Programmieren.
    \item \textbf{Übersetzen für alle Zielgeräte:} Die Zielgeräte müssen
    beim Kompilieren angegeben werden (siehe 4.1).
    \item \textbf{Fehlermeldungen:} Fehler im Kernel (z.\,B.\ zu große
    Work-Groups) zeigen sich erst zur Laufzeit und nur auf dem betroffenen
    Gerät.
\end{itemize}

\section{Grenzen der bisherigen Messungen}

\begin{itemize}
    \item Getestet wurden zwei Geräte (eine CPU, eine NVIDIA-GPU). AMD-
    und Intel-GPUs fehlen noch.
    \item Die Messreihen wurden jeweils einmal aufgenommen (mit fünf
    Wiederholungen pro Messpunkt). Kleine Unterschiede wie 1354 gegenüber
    1368\,GFLOPS können Messrauschen sein.
    \item Das Ausgangsprogramm (Stufe 0) ist absichtlich einfach gehalten.
    Die großen Beschleunigungsfaktoren gegenüber Stufe 0 sind deshalb kein
    fairer Vergleich mit optimiertem CPU-Code.
    \item In früheren Messungen (z.\,B.\ auf den Vortragsfolien) lief die
    CPU über AdaptiveCpp mit OpenMP statt über DPC++ mit OpenCL. Die
    Ergebnisse sind nicht direkt vergleichbar, weil der Compiler einen
    deutlichen Einfluss hat.
\end{itemize}

\section{Nächste Schritte und Fragen für das Gespräch}

\begin{enumerate}
    \item Messungen mehrfach wiederholen und Streuung angeben.
    \item Wenn möglich weitere Hardware testen (AMD- oder Intel-GPU).
    \item Vergleich der beiden SYCL-Compiler (DPC++ und AdaptiveCpp) auf
    derselben Hardware.
    \item Optional: Vergleich mit einer optimierten Bibliothek (z.\,B.\
    cuBLAS auf der GPU) als Obergrenze, um die Prozentwerte auf die
    tatsächlich erreichbare Leistung zu beziehen.
\end{enumerate}

\paragraph{Fragen an den Betreuer:}
\begin{itemize}
    \item Soll der Schwerpunkt eher auf dem \emph{Umstellungsaufwand} oder
    auf der \emph{Performance-Portabilität} liegen?
    \item Ist der Vergleich zweier Compiler im Rahmen der Seminararbeit
    gewünscht, oder reicht einer?
    \item Gibt es Zugang zu weiterer Hardware (z.\,B.\ AMD- oder
    Intel-GPU) über die Hochschule?
\end{itemize}

\end{document}
